
\label{sec:domain-eval}
\begin{table*}[]\small
\centering
\caption{Benchmarks for domain-specific factuality evaluation. The table presents the domain, tasks, datasets, and the LLMs evaluated in the respective study.}
\vspace{-2mm}
\label{tab:domain-eval}
\begin{adjustbox}{max width=\textwidth}
  \setlength{\tabcolsep}{2mm}
{
\begin{tabular}{lllll}
\toprule
Reference &
  Domain &
  Task &
  Metrics &
  %Human &
  Evaluated LLMs \\
  \midrule FLARE \citep{xie2023pixiu} &
  Finance  &
  \begin{tabular}[c]{@{}l@{}}Sentiment analysis, \\ News headline classification \\ Named entity recognition \\ Question answering \\ Stock movement prediction\end{tabular} &
  \begin{tabular}[c]{@{}l@{}}F1, ACC,\\Avg F1, \\ Entity F1, \\ EM, MCC\end{tabular} &
  %\checkmark &
  \begin{tabular}[c]{@{}l@{}}GPT-4 , \\ BloombergGPT, \\  FinMA-(7B, 30B, 7B-full), \\ Vicuna-7B \end{tabular}\\
   \midrule EcomInstruct \citep{li2023ecomgpt} & Finance  
   & 134 E-com tasks
   & \begin{tabular}[c]{@{}l@{}} Micro-F1, \\Macro-F1, \\
  ROUGE\end{tabular}
   & \begin{tabular}[c]{@{}l@{}}BLOOM, BLOOMZ,\\ ChatGPT, EcomGPT\end{tabular}
   \\
   \midrule CMB \citep{wang2023cmb} &
  Medicine &
  Multi-Choice QA &
  ACC &
  \begin{tabular}[c]{@{}l@{}}GPT-4, ChatGLM2-6B,\\ ChatGPT, DoctorGLM,\\ Baichuan-13B-chat,\\ HuatuoGPT, MedicalGPT,\\ ChatMed-Consult,\\ ChatGLM-Med ,\\ Bentsao, BianQue-2\end{tabular} \\
   \midrule Huatuo-26M \citep{li2023huatuo} &
Medicine & Generative-QA
    &
  \begin{tabular}[c]{@{}l@{}}BLEU, \\ROUGE, \\ GLEU\end{tabular} &
  T5, GPT2 \\
   \midrule NCBI \citep{jin2023genegpt} &
  Medicine &
  \begin{tabular}[c]{@{}l@{}}Nomenclature, \\ Genomic location, \\
  Functional analysis, \\
  Sequence alignment \end{tabular} &
  ACC & 
   \begin{tabular}[c]{@{}l@{}}GPT-2, BioGPT, \\
   BioMedLM, 
   GPT-3, \\
   ChatGPT, New Bing
   \end{tabular} \\
   \midrule LegalBench \citep{guha2023legalbench} &
  Law &
  \begin{tabular}[c]{@{}l@{}}Issue-spotting,  \\
  Rule-recall, \\
  Rule-application, \\ Rule-conclusion,
  \\
  Interpretation,
  \\
  Rhetorical-understanding
  \end{tabular}  &
  \begin{tabular}[c]{@{}l@{}} 
          ACC, EM
   \end{tabular}
  & 
   \begin{tabular}[c]{@{}l@{}}
   GPT-4, GPT-3.5, \\
   Claude-1, Incite, OPT\\
   Falcon, LLaMA-2, FLAN-T5...
   \end{tabular} \\
   \midrule 
   LawBench \citep{fei2023lawbench} &
  Law &
  \begin{tabular}[c]{@{}l@{}}Legal QA, NER, \\
  Sentiment Analysis, \\
  Reading Comprehension
  \end{tabular}  &
  \begin{tabular}[c]{@{}l@{}} 
         F1, ACC,\\
        ROUGE-L,\\
         Normalized \\log-distance	...
   \end{tabular}
  & 
   \begin{tabular}[c]{@{}l@{}}GPT-4,\\ 
   ChatGPT, \\
   InternLM-Chat,\\
   StableBeluga2...
   \end{tabular} \\
   \midrule 
   OceanBench \citep{bi2023oceangpt} &
  Geoscience &
  \begin{tabular}[c]{@{}l@{}}15  ocean-\\related tasks
  \end{tabular}  &
  \begin{tabular}[c]{@{}l@{}} 
         ACC
   \end{tabular}
  & 
   \begin{tabular}[c]{@{}l@{}}GPT-4
   \end{tabular} \\
  \bottomrule
\end{tabular}
}
\end{adjustbox}

\end{table*}

To assess the performance and factuality of these specialized LLMs, a plethora of datasets and benchmarks have been proposed across various domains. These resources not only serve as critical tools for evaluating the capabilities of LLMs but also facilitate advancements in specialized applications. We summarize them in Table~\ref{tab:domain-eval}. The distinction between this subsection and the previous two lies in its focus. This subsection delves deeper into factuality evaluation tailored to specific domains, while Sec \ref{sec:eval_benchmark} and \ref{sec:factuality-eval} primarily concentrate on general factuality evaluation, with only a portion of their content dedicated to datasets evaluating factuality within specific domains.


{\paragraph{Finance}}
\citet{xie2023pixiu} designed a financial natural language understanding and prediction evaluation benchmark dubbed FLARE, based on their collected financial instruction tuning dataset FIT. This benchmark is used to evaluate their FinMA model. It randomly selects validation sets from FIT to choose the best model checkpoint and utilizes distinct test sets for evaluation.
FLARE is a broader variant compared to the existing FLUE benchmark~\cite{shah2022flue} as it also encapsulates financial prediction tasks like stock movement prediction in addition to standard NLP tasks. 
The FLARE dataset includes several subtasks, such as sentiment analysis (FPB, FiQA-SA), news headline  classification (Headline), named entity recognition (NER), question answering (FinQA, ConvFinQA), and stock movement prediction (BigData22, ACL18, CIKM18).
Performance is gauged via a variety of metrics for each task, such as the accuracy and weighted F1 Score for sentiment analysis, entity-level F1 score for named entity recognition, and accuracy and Matthews correlation coefficient for stock movement prediction.
Several methods, including their own FIT-fine-tuned FinMA and other LLMs (BloombergGPT, GPT-4, ChatGPT, BLOOM, GPT-NeoX, OPT-66B, Vicuna-13B) are dedicated to their comparison. BloombergGPT's performance is assessed in various shot scenarios, while the zero-shot performance is reported for the remaining results. Some of the baselines depend on human evaluations given that LLMs without fine-tuning fail to generate instruction-defined answers. Conversely, FinMA's results are conducted on a zero-shot basis and can be evaluated automatically. 
To enable direct comparison between the performance of FinMA and BloombergGPT, despite the former not releasing their test datasets, test datasets were constructed with the same data distribution.

\citet{li2023ecomgpt} propose EcomInstruct benchmark, which experiment investigates the performance of their EcomGPT language model in comparison to baseline models such as BLOOM and BLOOMZ. Categories of these baseline models include pre-trained large models with decoder-only architecture like BLOOM, and instruction-following language models like BLOOMZ and ChatGPT. The evaluation metric of the EcomInstruct involves converting all tasks to generative paradigms and using text generation evaluation metrics like ROUGE-L. Classification tasks were evaluated with precision, recall, and F1 scores. The EcomInstruct dataset, comprising 12 tasks across four major categories, is divided into training and testing sections. The EcomGPT is trained on 85,746 instances of E-commerce data. The performance of the model is assessed based on its capacity to generalize unseen tasks or datasets, with emphasis on cross-language and cross-task paradigm settings.


{\paragraph{Medicine}}
\citet{wang2023cmb} propose a localized medical benchmark called CMB, or the Comprehensive Medical Benchmark in Chinese. CMB is rooted entirely in the native Chinese linguistic and cultural framework. While traditional Chinese medicine is a significant part of this evaluation, it does not make up the entire benchmark.
Both prominent LLMs, such as ChatGPT and GPT-4, and localized Chinese LLMs, including those specializing in the health domain, are evaluated using CMB. However, the benchmark is not designed as a leaderboard competition, but rather as a tool for self-assessment and understanding the progression of models in this field.
CMB embodies a comprehensive, multi-layered medical benchmark in Chinese, comprised of hundreds of thousands of multiple-choice questions and complex case consultation questions. This wide range of queries covers all clinical medical specialties and various professional levels, seeking to evaluate a model's medical knowledge and clinical consultation capabilities comprehensively.
 
\citet{li2023huatuo} introduce Huatuo-26M dataset, the largest Chinese medical Question and Answer (QA) dataset to date, including over 26 million high-quality medical QA pairs. It covers a wide range of topics, such as diseases, symptoms, treatments, and drug information. The dataset is a valuable resource for anyone looking to improve AI applications in the medical field, such as chatbots and intelligent diagnostic systems.
The Huatuo-26M dataset is gathered and integrated from various sources, including online medical encyclopedias, online medical knowledge bases, and online medical consultation records. Each QA pair in the dataset contains a problem description and a corresponding answer from a doctor or expert. Although a significant proportion of the Huatuo-26M dataset is constituted by the online medical consultation records, the text format data from these records is not publicly available, for unspecified reasons.
This dataset is expected to be instrumental for multiple types of research and AI applications in the medical field. These areas of application extend to Natural Language Processing tasks like developing QA systems, text classification, sentiment analysis, and Machine Learning model training tasks like disease prediction and personalized treatment recommendation. Consequently, the dataset is favorable for developing AI applications in the medical field, from intelligent diagnosis systems to medical consultation chatbots.

% {\paragraph{Gene: }}
\citet{jin2023genegpt} use nine tasks related to NCBI resources to evaluate the proposed GeneGPT model. The dataset used is GeneTuring benchmark~\cite{Hou2023GeneTuring}, which contains 12 tasks with 50 question-answer pairs each. The tasks are divided into four modules - Nomenclature  (Gene alias, Gene name conversion),  Genomic location(Gene SNP association, Gene location, SNP location), 
  Functional analysis(Gene disease association, Protein-coding genes), 
  Sequence alignment (DNA to human genome, DNA to multiple species). Two settings of GeneGPT are assessed: a full setting where all prompt components are used, and a slim setting using only two components. The performance of GeneGPT is compared against various baselines including general-domain GPT-based LLMs like GPT-2, GPT-3, and ChatGPT. Additionally, GPT-2-sized biomedical domain-specific LLMs such as BioGPT and BioMedLM are evaluated. The new Bing, a retrieval-augmented LLM with access to relevant web pages, is also assessed. The result evaluation of the compared methods is based on the results reported in the original benchmark and are manually evaluated. However, the evaluation of the proposed GeneGPT method is determined through automatic evaluations.





\paragraph{Law}
LegalBench~\citep{guha2023legalbench} is a benchmark for legal reasoning introduced due to the increasing use of LLMs in the legal field. It consists of 162 tasks covering six types of legal reasoning and was collaboratively constructed through an interdisciplinary process with significant contributions from legal professionals. The tasks designed aim to either demonstrate practical legal reasoning capabilities or measure reasoning skills that are of interest to lawyers. To facilitate discussions between different fields relating to LLMs in law, LegalBench tasks correspond to popular legal frameworks for describing legal reasoning, thus creating a shared language between lawyers and LLM developers. The paper not only describes LegalBench, but also presents an evaluation of 20 different open-source and commercial LLMs and highlights the types of research opportunities that LegalBench can provide.

LawBench~\citep{fei2023lawbench} is an evaluation framework dedicated to assessing the capabilities of LLMs in relation to legal tasks. The context-specificity and high-stakes nature of the law field make it crucial to have a clear grasp of LLMs' legal knowledge and their ability to execute legal tasks. 
LawBench dives deep into three cognitive evaluations of LLMs: the capability to memorize crucial legal details, understanding legal texts, and applying legal knowledge to resolve complicated legal problems. A total of 20 diverse tasks have been put together, covering five main task categories: single-label classification, multi-label classification, regression, extraction, and generation. 
Throughout the evaluation process, 51 LLMs were extensively tested under LawBench, compiling a spectrum of language models including 20 multilingual LLMs, 22 Chinese-oriented LLMs, and 9 law-specific LLMs. The results concluded that GPT-4 ranked as the superior LLM in the law domain, significantly surpassing its competitors. Despite the noted improvements when LLMs were fine-tuned on law-specific texts, the study acknowledged that there is still a long road ahead in achieving highly reliable LLMs for legal tasks.

\citet{bi2023oceangpt}  design OceanBench to evaluate the capabilities of LLMs for oceanography tasks.
OceanBench includes a total of 15 type of ocean-related tasks, such as question-answering and description tasks.
The samples in OceanBench are are automatically generated from the seed dataset and have undergone manual verification by experts.